\documentclass[10pt]{article}
\usepackage[margin=20mm]{geometry}
\usepackage{amsmath,amssymb,mathtools,fontspec,unicode-math,xurl,hyperref}
\IfFontExistsTF{DejaVu Serif}{\setmainfont{DejaVu Serif}}{\setmainfont{Latin Modern Roman}}
\setmathfont{latinmodern-math.otf}
\setlength{\parindent}{0pt}\setlength{\parskip}{6pt}\sloppy
\title{Technical review of
Trigonometric K-matrices for finite-dimensional representations of quantum affine algebras}
\author{AI-assisted technical review}
\date{}
\begin{document}
\maketitle
Abstract: This initial review concerns the 37-page preprint by Andrea Appel and Bart Vlaar, arXiv:2203.16503v5 (2 November 2025), associated with DOI 10.4171/JEMS/1686. Two local inaccuracies are confirmed in the preliminary root-data and grading descriptions. Explicit examples and corrections are given below. The inspected construction of spectral K-matrices, rational descent and generic restricted irreducibility does not rely on the inaccurate general descriptions in the ways disproved here. The assessment is Minor, with the limits of the proof checks recorded explicitly.\par
\section*{Version, hypotheses and proof coverage}
The exact source is the supplied arXiv v5, not an assumed copy of the journal article published online on 11 August 2025. The arXiv record identifies Remark 7.8.2 as an addition in v5. Its extension beyond Kirillov–Reshetikhin modules is therefore retained as a version-specific statement. All 37 PDF pages were read as extracted text; the selected evidence and main-argument pages were also inspected as images.\par
The algebraic setting is untwisted affine type over \(\mathbb F=\overline{\mathbb C(q)}\), with the alternative non-root-of-unity complex specialization noted by the authors. The finite-dimensional category \(\mathcal C\) is the type-1, level-zero category from Section 2.3. The QSP construction retains the generalized Satake conditions, the parameter restrictions \((\gamma,\sigma)\in\Gamma\times\Sigma\), and the admissible gauge requirements of Section 4.2. Theorem 4.2.1 is not read as allowing arbitrary twisting automorphisms.\par
The following central steps were checked. For Proposition 4.3.1(d), a fixed nonzero grading bounds the number of occurrences of a selected \(E_i\); between these occurrences, the other generators belong to a finite-type positive subalgebra and act nilpotently on the finite module. This supplies finiteness of each spectral coefficient without requiring the erroneous cone formula below. The identity \(\psi\Sigma_z=\Sigma_{1/z}\psi\) explains the target of the intertwiner in (4.1); its admissible-gauge hypothesis is essential.\par
For Lemma 5.1.1, rank of the finite linear intertwining system is unchanged on extending \(\mathbb F(z)\) to a Laurent-series field. Thus a nonzero formal solution yields a nonzero rational solution. Once the intertwiner space is one-dimensional, clearing denominators and removing the common polynomial factor gives the primitive polynomial representative in Theorem 5.1.2(b). The term non-vanishing here includes removal of common zeros, not merely exclusion of the identically zero matrix.\par
The support/depth induction in Theorem 6.2.4 was reconstructed: a monomial carrying a lower weight component to the top weight kills the top component, reducing support; otherwise an \(E_i\) decreases the chosen minimal depth. Distinct weight components remain distinct under a common weight shift, so cancellation does not invalidate this step. This uses Proposition 6.2.1 as an external input. For the specialization argument in Section 6.3, polynomial lifts of a basis of the fiber have determinant nonzero at zero and hence form a basis over \(\mathbb F(z)\). Under the principal grading, \(z\Sigma_z(\widetilde F_i)\) specializes to \(F_i\), since every correction lies in the nonnegative Borel. This verifies the deformation mechanism used for generic restricted irreducibility.\par
\section*{Confirmed local findings}
Problem 1 (Minor: incomplete affine positive-root descriptions).  On p. 6, Section 2.1 states \[\widehat\Delta^{\rm re}_+=\Delta_++\mathbb Z_{\ge0}\delta.\] On p. 12, Remark 3.2.1(c) also states \[\widehat Q_+=Q_+\oplus\mathbb Z_{\ge0}\delta.\] With the positive-root and nonnegative-cone conventions introduced on pp. 5–6, both right-hand sides omit the affine simple root.\par
Decisive example and assumptions. Take untwisted type \(A_1^{(1)}\), so \(\Delta_+=\{\alpha_1\}\), \(\vartheta=\alpha_1\), and \(\delta=\alpha_0+\alpha_1\). The affine simple root \(\alpha_0=\delta-\alpha_1\) is a positive real root and an element of \(\widehat Q_+\). Its coefficient of \(\alpha_1\), in the independent basis \((\alpha_1,\delta)\), is \(-1\). Every element of the two displayed right-hand sides has nonnegative finite-root coefficient (and the first has coefficient exactly \(1\)). Thus \(\alpha_0\) belongs to neither. For the setting of Remark 3.2.1, one may take the allowed generalized Satake diagram \(X=\varnothing,\tau=\mathrm{id}\); the root-cone contradiction is unchanged. There is no restriction on \(q\) or QSP parameters needed for this root-lattice calculation.\par
Equation reference: eq:cone. Correction and repair boundary. For the untwisted affine positive real roots, the missing branch is \[\widehat\Delta^{\rm re}_+
 =(\Delta_++\mathbb Z_{\ge0}\delta)
   \cup(-\Delta_++\mathbb Z_{>0}\delta).\] For the cone, the exact formula is \[\widehat Q_+
 =\{\,\beta+n\delta:\ \beta\in Q,\ n\in\mathbb Z_{\ge0},\
       \beta+n\vartheta\in Q_+\,\}.
 \] Indeed \(n\alpha_0+\gamma=(\gamma-n\vartheta)+n\delta\) for \(n\ge0,\gamma\in Q_+\), and this identity proves both directions of [eq:cone] (\#eq:cone). The group decomposition \(\widehat Q=Q\oplus\mathbb Z\delta\) does not yield a product decomposition of the positive cone.\par
Impact and counterevidence. These are false preliminary descriptions, but the checked descent proof uses monomials in all affine simple generators, and Section 6 explicitly uses the projection \(\overline{\alpha_0}=-\vartheta\) (p. 26). In particular, that proof does not assume that every positive affine root has positive finite projection. The rank-one implication \(\widehat Q^{-\theta}=\mathbb Z\delta\) can be obtained from the rank and the primitive element \(\delta\), without the false cone decomposition. No failure of Theorems 4.2.1, 5.1.2 or 6.1.1 has been established from these descriptions.\par
Problem 2 (Minor: invariant integer gradings need not be indicator functions).  On p. 16, Section 4.1 first allows any group homomorphism \(s:\widehat Q\to\mathbb Z\), then asserts that \(\tau\)-invariance is equivalent to \(s\), on affine simple roots, being the characteristic function of a union of \(\tau\)-orbits. The necessity assertion requires an additional \(0/1\)-valued restriction.\par
Decisive example. Again use \(A_1^{(1)}\), \(X=\varnothing,\tau=\mathrm{id}\), and define \(s(\alpha_0)=s(\alpha_1)=2\). Since the simple roots freely generate \(\widehat Q\), this is a group homomorphism satisfying all hypotheses of the general sentence. It commutes with the identity diagram automorphism and has \(s(\delta)=4>0\), but it is not an indicator function. Directly on generators, invariance is equivalent to \[s(\alpha_{\tau(i)})=s(\alpha_i)\quad\text{for every }i.\] Thus the correct general description is an arbitrary integer constant on each orbit. Equivalently, \(s=\sum_O m_O{\bf1}_O\) on simple roots, \(m_O\in\mathbb Z\).\par
Impact and repair boundary. Replace the indicator-function claim by orbitwise constancy, or explicitly restrict that sentence to \(0/1\)-valued gradings. The subsequently defined \(\tau\)-minimal grading is indeed \(0/1\)-valued and invariant; the principal grading used in Section 6 also satisfies the required invariance. The counterexample does not invalidate those selected gradings or the resulting spectral operators. The finding is local for this reason, not because the textual correction is short.\par
\section*{Boundary checks and limitations}
Propositions 5.4.1–5.5.1 retain the stated \(\psi\)-involutivity or \(\psi\)-fixedness assumptions and the irreducibility condition at a special value. The pole argument uses an invertible leading intertwiner and the non-vanishing normalization of its partner; it does not assert that reducibility must always produce a pole. The explicit constant representative in Section 7.9(c) has \[J_\nu=\begin{pmatrix}0&-\nu^{-1}\\ \nu&0\end{pmatrix},
 \qquad J_\nu^2=-I.\] Accordingly \(iJ_\nu\) supplies a unitary scalar normalization. This computation is consistent with a pole-free representative and with the preceding qualification that the solution is determined up to a scalar; it is not used to refute the existence theorem.\par
The shifts, amenability assumptions and the exceptional affine \(D_{2n}\) case in Section 7 were retained. The full universal K-matrix construction from the cited earlier work, the classification of generalized Satake diagrams, all R-matrix normalization identities, and every representation-theoretic identification in Section 7 were not independently reproved. Reading the whole body does not remove these dependencies.\par
The source manifest and bounded searches of the arXiv, publisher and author pages did not identify separate computational supplements or an author repository. No author code was executed; absence of code is not a mathematical defect. All source inspection, dependency assessment, reporting, compilation and page review were performed in the present initial review conversation. The scientific decision remains pending; no human review or independent acceptance is recorded.\par
9 A. Appel and B. Vlaar, Trigonometric K-matrices for finite-dimensional representations of quantum affine algebras, arXiv:2203.16503v5, 2 November 2025. https://arxiv.org/abs/2203.16503v5. Related journal DOI: 10.4171/JEMS/1686 (https://doi.org/10.4171/JEMS/1686).\par
\end{document}
